% BOOK-FILE-ID: f-8b932adc7bc1
ここまでは,式変形によって基本パターンに帰着させることを考えてきた.
しかし,世の中にはそれではどうにもならない漸化式も多い.
まずは次の3つの\bookterm{recurrence}{漸化式}を見てほしい.
\[
  a_{n+1}=a_n^2,
  \quad
  a_{n+1}=a_n^2-1,
  \quad
  a_{n+1}=a_n^2-2
\]
1つ目は簡単である.1回進むたびに指数が2倍になるので,書き下していけば
\[
  a_n=a_1^{2^{n-1}}
\]
と分かる.\,$a_1>0$なら\bookterm{logarithm}{対数}を取って等比型にもできる.\footnote{$a_1\le0$では実数の対数を最初から取れないが,元の漸化式は問題なく定義される.}

2つ目はどうだろうか.\,$-1$が邪魔をして,対数を取っても和の形にならない.
ならばと特性方程式型のように\bookterm{fixed}{不動点}を求めてみると
\[
  \alpha=\alpha^2-1
  \quad\Longleftrightarrow\quad
  \alpha=\frac{1\pm\sqrt5}{2}
\]
であり,初項をこれに選べば定数解となる(\sectionref{節：初項による違い}).
ただし,\,$b_n=a_n-\alpha$と置いても,
\[
  b_{n+1}=b_n^2+2\alpha b_n
\]
と2次の項が残ってしまい,\bookterm{linear}{線形}にはならない.
この漸化式については,本書で使う関数を用い,\,$n$が大きくなっても計算回数の増えない一般項を求める方法を紹介できない.\footnote{どのような式を答えとするか,また「解く」ことの範囲は\sectionref{節：解法の分類}で考える.}

では3つ目はどうか.2つ目とは定数項が$-1$か$-2$かの違いしかないので,同じく解けないと思うかもしれない.
ところが,この$-2$はある恒等式にちょうど合うのだ.
基本パターンへの変形が見つからなくても,その式だけが持つ特別な構造に気付けば,道が開けることがある.
この節では,そうした構造を使う解法を見ていこう.
ここで大事なのは,やはり実験と「この形は見たことがある」という勘である.
